\documentclass[
  spanish,
  letterpaper, oneside
]{article}

\def\documenttitle {El Estado a través de la historia}
\def\documentsubtitle {Evolución del poder y protección de los derechos humanos}
\def\documentsubject {Actividad 2 - Garantías Constitucionales}

\def\documentauthor {Martin Jonathan de la Cruz}
\def\coursename {Garantías Constitucionales}
\def\coursecode {025}

\def\universityname {Universidad Abierta y a Distancia de México}
\def\universityfaculty {Licenciatura en Derecho}
\def\universitydepartment {Garantías Constitucionales}
\def\universitydepartmentimage {departamentos/UnADM}
\def\universitydepartmentimagecfg {height=1.57cm}
\def\universitylocation {Roma Norte, Ciudad de México}

\def\authortable {
  \begin{tabular}{ll}
    \textbf{Alumno:}
    & \begin{tabular}[t]{l}
      Martin Jonathan de la Cruz \\
    \end{tabular} \\
    Matrícula: & ES2611202040 \\
    Figura docente: & Aarón López Pérez \\
    Semestre/Bloque: & 2 / 1 \\
    & \\
    \multicolumn{2}{l}{Fecha de realización: \today} \\
    \multicolumn{2}{l}{\universitylocation}
  \end{tabular}
}

\input{template}
\usepackage{helvet}
\renewcommand{\familydefault}{\sfdefault}
\usepackage{tikz}
\usetikzlibrary{calc}
\setcitestyle{authoryear,open={(},close={)}}

\def\coverwatermarkenabled {true}
\def\coverwatermarkimage {img/departamentos/UnADM.pdf}
\def\coverwatermarkopacity {0.16}
\def\coverwatermarkwidth {0.72\paperwidth}
\def\coverwatermarkxshift {0cm}
\def\coverwatermarkyshift {-0.35cm}

\newcommand{\insertcoverwatermark}{%
  \ifthenelse{\equal{\coverwatermarkenabled}{true}}{%
    \AddToShipoutPictureBG*{%
      \begin{tikzpicture}[remember picture,overlay]
        \node[
          opacity=\coverwatermarkopacity,
          inner sep=0pt,
          xshift=\coverwatermarkxshift,
          yshift=\coverwatermarkyshift
        ] at (current page.center) {%
          \includegraphics[width=\coverwatermarkwidth]{\coverwatermarkimage}%
        };
      \end{tikzpicture}%
    }%
  }{}%
}

\begin{document}

\insertcoverwatermark
\templatePortrait
\templatePagecfg
\onehalfspacing

\begin{abstractd}
  La evolución histórica del Estado muestra una transformación progresiva de las formas de organización y legitimación del poder. Desde las agrupaciones humanas y la ciudad-Estado hasta el Estado social y democrático de derecho, cada etapa modificó la relación entre autoridad, comunidad y libertad. El recorrido permite comprender que los derechos humanos y sus garantías no aparecieron de manera aislada, sino como respuesta a experiencias históricas de desigualdad, concentración del poder y exclusión.
\end{abstractd}

\templateIndex
\templateFinalcfg

\section{Introducción}

El Estado no es una institución inmóvil. Sus características actuales son resultado de procesos históricos en los que las comunidades organizaron el poder, delimitaron territorios, establecieron autoridades y construyeron reglas para resolver conflictos. La nota local de la asignatura identifica una secuencia que va de las agrupaciones humanas prehistóricas y la ciudad-Estado al feudalismo, el absolutismo, el Estado moderno, el liberalismo y el Estado social y democrático de derecho \citep{notaEstadoHistoria2026}.

\textbf{Objetivo:} analizar la evolución histórica del Estado para explicar cómo la limitación jurídica del poder favoreció el reconocimiento y la protección de los derechos humanos.

\textbf{Propósito de la participación:} sostener una postura personal fundamentada, relacionar las etapas históricas con las garantías constitucionales y abrir el intercambio con una pregunta que permita contrastar argumentos.

\textbf{Criterios de evaluación:} claridad del recorrido histórico, uso pertinente de fuentes locales, relación entre poder y derechos humanos, argumentación propia y formulación de una pregunta relevante para el diálogo.

Mi postura es que esta evolución puede leerse como una tensión constante entre concentración y limitación del poder. Las formas antiguas y medievales permitieron organizar la vida colectiva, pero también mostraron que una autoridad sin límites suficientes puede convertir a las personas en súbditos sin medios efectivos de defensa. La construcción histórica de los derechos humanos responde precisamente a la necesidad de reconocer la dignidad y establecer límites jurídicos frente al poder \citep{rabinovichBerkmanDerechosHumanos,marcanoSalazarDerechosHumanos}.

\section{Participación en el foro}

\subsection{De las agrupaciones humanas a la ciudad-Estado}

Las primeras agrupaciones humanas surgieron de necesidades compartidas de subsistencia, seguridad y cooperación. Con el tiempo, formas como la horda, el clan, la fratría y la tribu hicieron posible una identidad colectiva más estable. La idea aristotélica del ser humano como \textit{zoon politikon} resume que la vida política no es un elemento accidental, sino una dimensión de la convivencia humana.

La polis griega y la civitas romana representaron un avance porque vincularon territorio, ciudadanía, normas y participación pública. Sin embargo, la condición de ciudadano no era universal. La exclusión de mujeres, personas esclavizadas y extranjeros revela que la participación política coexistía con profundas desigualdades. Esta contradicción ayuda a entender por qué la historia de los derechos humanos no consiste únicamente en reconocer instituciones políticas, sino en ampliar progresivamente quién puede ser considerado titular de derechos \citep{rabinovichBerkmanDerechosHumanos}.

\subsection{Feudalismo y absolutismo: poder fragmentado y poder concentrado}

En el Estado feudal el poder se fragmentó entre monarquías, nobleza, señores territoriales e Iglesia. Las relaciones de vasallaje y dependencia personal sustituyeron una autoridad pública centralizada. La posición jurídica de las personas dependía en gran medida del estamento al que pertenecían, por lo que no existía igualdad general ante la ley.

El absolutismo respondió a esa fragmentación mediante la centralización territorial, administrativa y militar. La soberanía permitió explicar la supremacía del poder estatal, pero su identificación con la voluntad del monarca produjo un nuevo problema: el poder unificado podía actuar sin controles suficientes. La importancia de esta etapa radica en que hizo visible la necesidad de separar la existencia del Estado de la voluntad personal de quien gobierna.

\subsection{Estado moderno y Estado liberal: límites jurídicos al poder}

Las ideas de soberanía popular, división de poderes, libertad e igualdad jurídica transformaron la legitimidad estatal. El poder dejó de justificarse exclusivamente por herencia o voluntad divina y comenzó a relacionarse con la nación, la ciudadanía y la Constitución. Las revoluciones modernas impulsaron declaraciones de derechos que reconocieron libertades civiles y políticas, aunque su aplicación inicial siguió siendo incompleta y excluyente.

El Estado liberal aportó un principio decisivo para las garantías constitucionales: la autoridad debe estar sometida al derecho. No basta con que exista un gobierno organizado; también se requieren competencias delimitadas, división de poderes y mecanismos para cuestionar los actos que vulneren libertades. Las teorías contemporáneas de los derechos humanos refuerzan esta relación entre dignidad, limitación del poder y exigibilidad jurídica \citep{marcanoSalazarDerechosHumanos}.

\subsection{Estado social y democrático de derecho}

La igualdad formal del liberalismo resultó insuficiente frente a condiciones materiales de pobreza, explotación y exclusión. Por ello, el Estado social amplió sus responsabilidades hacia la educación, la salud, el trabajo, la seguridad social y otras condiciones necesarias para una vida digna. El Estado democrático de derecho integra participación política, legalidad, control del poder y protección de derechos civiles, políticos, económicos, sociales y culturales.

En México, el reto no consiste únicamente en reconocer derechos en el orden jurídico, sino en lograr que las instituciones los hagan efectivos para todas las personas. La distancia entre reconocimiento formal y experiencia cotidiana explica la importancia de organismos, procedimientos y garantías que permitan prevenir, denunciar y reparar violaciones de derechos humanos \citep{alvarezIcazaDerechosMexico}.

\subsection{Relación con las garantías constitucionales}

El recorrido histórico demuestra que los derechos necesitan instrumentos de protección. Una declaración puede formular libertades e igualdad, pero pierde eficacia si la autoridad no está obligada a respetarla o si las personas carecen de vías para reclamar. Las garantías constitucionales convierten los límites del poder en procedimientos e instituciones exigibles.

Considero que el principal aprendizaje es distinguir entre la existencia del Estado y su legitimidad. Un Estado puede disponer de territorio, población, gobierno y soberanía, pero solo se aproxima a un Estado constitucional cuando el ejercicio del poder se encuentra jurídicamente limitado, existe participación democrática y las personas cuentan con medios efectivos para defender su dignidad y sus derechos.

\section{Pregunta para el intercambio}

¿Puede considerarse democrático un Estado que reconoce una amplia lista de derechos en su Constitución, pero no ofrece garantías accesibles y eficaces para hacerlos valer? En mi opinión, no es suficiente el reconocimiento formal: la calidad constitucional del Estado también debe evaluarse por la posibilidad real de prevenir abusos, obtener protección y reparar violaciones.

\section{Conclusión}

La historia del Estado muestra un tránsito desde organizaciones basadas en parentesco, estamento o concentración personal del poder hacia modelos que buscan fundamentarse en ciudadanía, legalidad, democracia y derechos humanos. Este tránsito no fue lineal ni definitivo; cada avance coexistió con exclusiones y conflictos que obligaron a ampliar el significado de libertad, igualdad y dignidad.

Mi postura final es que el Estado social y democrático de derecho representa un ideal exigente, no una condición automática. Su legitimidad depende de que la Constitución limite efectivamente a las autoridades y de que las garantías permitan convertir los derechos en protección concreta. La consecuencia práctica es clara: estudiar garantías constitucionales implica analizar no solo qué derechos se reconocen, sino también quién puede ejercerlos, frente a quién y mediante qué mecanismos.

\clearpage
\nocite{notaEstadoHistoria2026,rabinovichBerkmanDerechosHumanos,marcanoSalazarDerechosHumanos,alvarezIcazaDerechosMexico}
\bibliography{garantias-constitucionales}

\end{document}